\section{Differentiability and the boundary}
\label{df:sec:conclusion}

The preceding section has ruled out distinct extreme supports at
every prescribed direction. We now pass from this directional
statement to Fr\'echet differentiability and identify its equivalent
boundary formulation.

\begin{lemma}[Convex boundary regularity]\label{df:convex-regularity}
For a norm on $\R^2$, Fr\'echet differentiability at every nonzero
vector is equivalent to uniqueness of the supporting line at every
point of its unit sphere. These conditions are also equivalent to
the unit sphere being a regular $C^1$ curve.
\end{lemma}

\begin{proof}
\emph{Differentiability and a unique support.}
Fix a nonzero vector $u$ and choose a vector $e$ linearly independent
of it. The function $t\mapsto\mu(u+te)$ has finite one-sided
derivatives $a_-\le a_+$. As in Lemma~\ref{df:prelim:supports}, the
linear functionals taking the value $\mu(u)$ on $u$ and the values
$a_-$ and $a_+$ on $e$ support $\mu$ globally. Any other supporting
functional at $u$ must have its value on $e$ between these two
numbers, by the inequalities for positive and negative $t$.

If the endpoints agree with value $a$, the two-sided expansion gives,
as $xu+re\to u$ with $x>0$,
\[
 \mu(xu+re)
 =x\mu\bigl(u+(r/x)e\bigr)
 =x\mu(u)+ar+o(|r|).
\]
Since $x\to1$ and $|r|\le C|xu+re-u|$, this is Fr\'echet
differentiability at $u$. The expression is exact when $r=0$.
Conversely, Fr\'echet differentiability makes the two one-sided
derivatives equal. Thus differentiability is equivalent to a unique
linear functional $g\le\mu$ satisfying $g(u)=\mu(u)$.

\emph{Supports of the norm and of its unit ball.}
Because zero lies in the interior of the unit ball, every supporting
line at $u/\mu(u)$ has a unique normalization $g=1$ with
$g\le\mu$ and $g(u)=\mu(u)$. This identifies uniqueness of the
functional with uniqueness of the line.

\emph{Continuity of the derivative and the boundary.}
Suppose now that this uniqueness holds everywhere off zero. All
supporting functionals lie in the compact dual unit ball. If
$u_j\to u\ne0$ and a subsequence of the supports at $u_j$ converges
to $g$, passing to the limit shows $g\le\mu$ and $g(u)=\mu(u)$.
Uniqueness determines $g$, so the whole sequence of supports
converges to the support at $u$. The derivative of $\mu$ is therefore
continuous off zero. At a unit vector for this norm it is nonzero,
since homogeneity gives
\[
 D\mu(u)[u]=\mu(u)=1.
\]
The implicit function theorem shows that the level set $\mu=1$ is a
$C^1$ curve.

Conversely, at any point of a regular $C^1$ convex boundary a supporting
functional has a local maximum along a local boundary
parametrization. It therefore annihilates the tangent direction.
Its supporting line is the tangent line, which is unique. This
proves the converse and all the equivalences.
\end{proof}

\begin{proof}[Proof of Theorem~\ref{thm:gamma-differentiability}]
Proposition~\ref{df:wi:no-corner} gives $a_-=a_+$ for every fixed
$b\in[-1,1]$. The expansion in the proof of
Lemma~\ref{df:convex-regularity}, applied to $u=(1,b)$ and $e=e_2$,
gives Fr\'echet differentiability at $u$. Signed coordinate
permutations and positive homogeneity cover every nonzero vector.
Lemma~\ref{df:convex-regularity} gives the remaining conclusions.
The parameters $\kappa,\lambda>0$ were arbitrary and fixed throughout,
so the conclusion holds for each Gamma law in the theorem.
\end{proof}
